\section{三相と分母：何に対する割合か}\label{theory:engineering:phases}
土粒子の体積を\(V_s\)、水を\(V_w\)、空気を\(V_a\)とし、
全体積を\(V=V_s+V_w+V_a\)、間隙体積を\(V_v=V_w+V_a\)とする。
乾燥土粒子の質量を\(M_s\)、水質量を\(M_w\)とすれば
\[
 w=\frac{M_w}{M_s},\quad n=\frac{V_v}{V},\quad
 e=\frac{V_v}{V_s},\quad \theta=\frac{V_w}{V},\quad
 S_r=\frac{V_w}{V_v}.
\]
順に含水比、間隙率、間隙比、体積含水率、飽和度である。
百分率で表す場合は該当する比に100を掛ける。\(n=e/(1+e)\)、\(\theta=nS_r\)は
同じ三相の体積から導ける。含水比は水が乾燥土より重ければ1を超えるが、
体積の定義に基づく\(n,\theta,S_r\)は通常0から1である。\cite{eng52_soil_phases}

湿潤質量\(M_s+M_w\)で割る含水率と、乾燥質量で割る土質の含水比を取り違えない。
質量を体積で割れば密度の次元が残るので、百分率の候補を除外できる。
乾燥質量は測定条件に依存する。平成25年度の指定\(110\pm5\,^\circ\mathrm C\)を保持し、
別の試験規約の温度へ勝手に置き換えない。

\section{有効応力から液状化を説明する}\label{theory:engineering:liquefaction}
飽和土の骨格へ作用する応力を、説明用に\(\sigma'=\sigma-u\)と表す。
全応力\(\sigma\)と間隙水圧\(u\)を分ける。緩い砂質土が短時間の繰返しせん断を受け、
水が十分に排出できない場合、収縮しようとする骨格の応答が過剰間隙水圧へ現れる。
uが上がると有効応力が下がり、抵抗・剛性が低下して大きな変形が生じ得る。\cite{eng52_liquefaction}

高い地下水位、若く緩い砂質堆積物、よく淘汰された粒度等は代表的な条件である。
旧河道・埋立地・浜堤列間の低地という地形履歴を地下水・粒度・密度へつなぐ。
N値は標準貫入の抵抗の指標で、低いNは緩い砂の傾向を示すが、上載圧や細粒分・地震外力も考える。
原典のN20以上やN10以下を全土層の安全・危険を分ける一つの絶対線とはしない。
淘汰の悪い土も組成・密度によって挙動が異なる。\cite{eng52_liq_sorting}

試料を採取して動的な変形・液状化抵抗を調べ、PS検層で応答に関係する波速を得る。
試料の乱れと現場の密度を確認する。通常の平板載荷は静的な支持・沈下が主目的なので、
発生の評価の相対比較では直接関係が薄い。しかし液状化後の地盤反力を平板載荷で測る研究もある。
「最も関係が薄い」を「絶対に利用されない」へ変えない。\cite{eng52_plate_liq,geebec_ps}

\section{試験名を測定量と条件へ戻す}\label{theory:engineering:tests}
三軸試験は拘束圧と排水・圧密の条件を指定し、破壊と変形を測る。
一軸試験は側方拘束を与えずに圧縮する。ベーン試験は回転のせん断抵抗、
標準貫入やサウンディングは貫入抵抗を調べる。
前二者の破壊試験と、貫入から経験的に推定する強度・支持性を同じ直接測定とはしない。\cite{eng52_geotech_tests}
圧密試験は荷重の増加に対する時間依存の体積変化・圧縮性・排水を調べる。
状態が変われば強度も変わり得るが、主目的はせん断破壊強度の直接測定ではない。\cite{eng52_consolidation}

\section{Mohr--Coulombの切片・傾きと測定丸め}\label{theory:engineering:mohr}
原典が直線の包絡線を指定したときは
\[
 \tau_f=c+\sigma_n\tan\phi
\]
として切片cと傾き\(\tan\phi\)を分ける。
有効応力の強度定数を使う場合は
\(\tau_f=c'+(\sigma_n-u)\tan\phi'\)となるが、問題の記号を無断で有効応力へ変更しない。
摩擦角の値をそのまま傾きへ入れず、正接を使う。\cite{eng52_mohr_ref}

平成29年度の測定は共通単位\(10^4\,\mathrm{Pa}\)で
\((1,3.1),(2,3.7),(4,4.8),(8,7.1)\)。
端点の傾きは\(4/7\simeq0.5714\)で指定\(\tan30^\circ=0.58\)に近い。
各点の\(\tau-0.58\sigma\)は2.52、2.54、2.48、2.46で、切片2.5の候補と整合する。
丸めた測定点は厳密には一直線ではない。通常の最小二乗でも
傾き0.5696、切片2.5391となり同じ選択となる。数値単位はcにも掛かる。
残差をゼロにする架空の測定値へ修正しない。

\section{揚水と圧密：水が抜けて骨格へ荷重が移る}\label{theory:engineering:consolidation}
地下水揚水で被圧帯水層の水頭が下がると、隣接する粘土層の水圧にも影響する。
同じ全応力でuが低下すれば有効応力は増加する。透水性が低い粘土では、
排水に時間を要して収縮し、地表の沈下として現れる。\cite{eng52_consolidation,eng52_subsidence_ref}
砂の即時的な変形、粘土の圧密、地震による液状化後の体積変化を区別する。
水頭低下と沈下が同時同量でないことにも注意する。

杭等に支えられた構造物の相対的な浮き上がり、道路・埋設物の変形に加え、
低平地では排水・浸水リスクも変わる。深度ごとに帯水層の水頭を観測し、
複数層をつないだ混合水位だけで原因を判断しない。
不圧地下水位の低下だけを広域沈下の唯一の機構としない。

\section{黄鉄鉱の酸化と、溶解・析出の時間順}\label{theory:engineering:pyrite}
還元的な環境で保存された黄鉄鉱は、掘削後に水と酸素へさらされると酸化する。
例えば二価鉄を生成する段階を収支式で書くと
\[
 \mathrm{FeS_2}+\frac72\mathrm{O_2}+\mathrm{H_2O}
 \longrightarrow \mathrm{Fe^{2+}}+2\mathrm{SO_4^{2-}}+2\mathrm{H^+}.
\]
さらに鉄の酸化・加水分解、炭酸塩による中和等が加わるので、実際のpHは材料と水の条件に依存する。
酸化そのものをアルカリ性排水の起源とはしない。\cite{eng52_weathering}

酸性水が膠結物等を溶かして岩石を弱める過程と、溶けた成分が石膏等の硫酸塩で再析出する過程を分ける。
蒸発・過飽和・水の経路によって析出し、結晶成長圧が盤膨れへつながる例がある。
青黒い未酸化岩の性質を、掘削後も酸素に触れないまま硫黄の酸化が進むという説明へ変えない。
酸化帯の充填、下の溶解帯の弱化などがあるため、酸化の存在だけで強化又は抑制を一方向へ決めない。\cite{eng52_weathering}

\section{鉄の色・還元溶出と交換サイト}\label{theory:engineering:soilredox}
酸化的な鉄の酸化物等は赤褐色・黄褐色、還元的なグライ化は青灰色等という代表的な対応がある。
土壌の色は腐植等にも依存するので一つの鉄化合物だけから全部を決めない。
湛水して酸素が消費されるとFe・Mnが還元され、二価の形で溶出しやすくなる場合がある。\cite{eng52_soil_colour}
水圧による力学的な不安定化と、酸化還元による化学的な物性変化を分ける。

粘土・腐植の交換サイトではカチオンが電荷を保ちながら入れ替わる。
同じ固定サイト系での保存条件と、pH等に伴う材料の交換容量の変化を区別する。
土壌有機炭素は重要だが、炭酸塩に富む土壌を含む全炭素を有機物だけへ限定しない。

\section{斜面：弱面と地下水を幾何へ重ねる}\label{theory:engineering:landslide}
重力性の移動を物質・移動様式・速度の履歴で記述する。
層理面・変質凝灰岩等の弱層が斜面へ抜ける流れ盤ではすべりへ注意する。
背斜・ドーム・キャップロックという名前だけで全斜面の安定を決めず、割れ目・水・地形を重ねる。\cite{eng52_slip_weaklayer}
すべり面の抵抗を\(c'+(\sigma_n-u)\tan\phi'\)で考えると、
水圧の上昇で抵抗が下がり得ることを説明できる。

滑落崖、亀裂、隆起・陥没は移動を示す観察対象である。長期緩慢・反復する地すべりも、
急激な滑動もある。「地すべり＝全例低速、崩壊＝全例脆性」の二分法は使わない。
原典の砂質/粘性の逆転は典型比較の中で説明し、全岩種を排他的に分類する根拠にはしない。\cite{eng52_landslide_ref}
浸透を減らす被覆・地下水排水・押え盛土等は機構に応じた対策候補であり、
現地の安全と施工条件を確認せず実行する一般指示ではない。

\section{トンネル：切羽・緩み・応力集中}\label{theory:engineering:tunnel}
掘削は空洞周囲の応力を変え、不連続面に沿う落下・緩み、水による切羽の弱化等を生じ得る。
切羽の天端・鏡面を観察し、地質構造と湧水を支保・設計へ反映する。\cite{eng52_tunnel_ref}
山はねは代表的には高応力・大きな土かぶり・硬く節理の少ない岩盤で問題になり、
小土かぶり・多節理を典型条件とする原典の候補は逆である。\cite{eng52_rockburst}
斜面地形・傾斜地層等の非対称な偏圧、浅部の地表沈下、地下水低下による沈下を分ける。
深さだけで沈下量を一つの単調則として決めない。

\section{骨材：岩種は品質試験の入口}\label{theory:engineering:aggregate}
岩石名から鉱物組成・強度・耐久性・吸水・反応性を考えるが、名称だけで品質を保証しない。
蛇紋岩骨材のポップアウト研究はある一方、骨材用採石の記録もある。
公式の一般的な不適性の比較と、全試料の使用不可能という全称を区別する。\cite{eng52_serpent_aggregate,eng52_popout}
チャートのシリカの反応性、砂岩の結合、片麻岩の方向性等も考える。
一つの健全試料の結果を全変質・風化状態へ広げない。

\section{浅部地中熱：伝導と移流を分離する}\label{theory:engineering:groundheat}
浅部の安定した温度と外気との差を利用して冷暖房する地中熱を、深部の高温地熱発電と区別する。
閉じた地中配管で媒体を循環するクローズドループと、地下水を利用するオープンループを分ける。
地盤・帯水層を熱の貯蔵へ使う場合は蓄熱の履歴と周辺環境への影響も考える。\cite{eng52_groundheat_ref}

熱伝導は温度勾配による輸送、地下水の移流は水が熱を運ぶ過程である。
熱応答試験を伝導モデルで整理すると移流の効果が見掛けの有効熱伝導率へ現れ、
材料自体の伝導率より大きく推定される例がある。\cite{eng52_heat_advection}
原典の一般的な大幅減少という向きを選ぶ誤りとして区別する。
その寄与は流速・方向・幾何・試験時間に依存し、全試験の一定倍率ではない。
